\documentclass[11pt,twocolumn]{article}

%% Reasoning-aware Adaptive Routing for Hybrid GraphRAG
%% Authors: Cong Tran, Nguyen Dinh Thanh, Pham Van Hai

% ============================================================
% Encoding and font
% ============================================================
\usepackage[utf8]{inputenc}
\usepackage[T5]{fontenc}
\usepackage[vietnamese,english]{babel}
\selectlanguage{english}
\usepackage{lmodern}

% ============================================================
% Math
% ============================================================
\usepackage{amsmath,amssymb,amsthm,mathtools}

% ============================================================
% Tables
% ============================================================
\usepackage{booktabs,tabularx,multirow,array,makecell,adjustbox}
\usepackage{ragged2e}
\newcolumntype{Y}{>{\RaggedRight\arraybackslash}X}
\newcolumntype{L}[1]{>{\RaggedRight\arraybackslash}p{#1}}

% ============================================================
% Figures, algorithms, layout
% ============================================================
\usepackage{graphicx}
\usepackage{xcolor}
\usepackage{algorithm}
\usepackage{algorithmicx}
\usepackage{algpseudocode}
\usepackage{float}
\usepackage{placeins}
\usepackage{titlesec}
\raggedbottom
\setlength{\parskip}{0pt}
\setlength{\parindent}{1em}
\titlespacing*{\section}{0pt}{1.0ex plus 0.3ex minus 0.2ex}{0.7ex plus 0.2ex}
\titlespacing*{\subsection}{0pt}{0.8ex plus 0.25ex minus 0.2ex}{0.45ex plus 0.15ex}
\titlespacing*{\subsubsection}{0pt}{0.7ex plus 0.2ex minus 0.2ex}{0.35ex plus 0.1ex}

% ============================================================
% pgfplots (for inline bar charts)
% ============================================================
\usepackage{pgfplots}
\pgfplotsset{compat=1.18}

% ============================================================
% Hyperref (last)
% ============================================================
\usepackage{hyperref}

% ============================================================
% Suppress overfull hbox warnings
% ============================================================
\hbadness=10000
\vbadness=10000

% ============================================================
% Abstract environment (Springer spmpsci style)
% ============================================================
\newtheorem{proposition}{Proposition}
\def\abstract{\typeout{Abstract}{\bf Abstract}\par}
\def\endabstract{\par}

% ============================================================
\begin{document}

\title{Reasoning-aware Adaptive Routing for Hybrid GraphRAG}

\author{
  Cong Tran$^1$ \and
  Nguyen Dinh Thanh$^2$ \and
  Pham Van Hai$^2$
}
\date{
  $^1$Posts and Telecommunications Institute of Technology (PTIT), Hanoi, Vietnam\\
  $^2$School of Information Technology and Communication,
  Hanoi University of Science and Technology, Hanoi, Vietnam\\
  \texttt{thanh.nd225670@sis.hust.edu.vn}
}

\maketitle

% ============================================================
\begin{abstract}
In Vietnamese legal question answering, the appropriate retrieval
strategy varies with the query: simple questions require direct textual
lookup, relational questions demand reasoning over legal structures such
as amendment, replacement, and cross-document reference, while ambiguous
questions risk grounding answers in the wrong legal target. We propose a
Hybrid GraphRAG system with a two-stage Reasoning-aware Adaptive Router
that selects among dense retrieval, graph traversal, hybrid reasoning, and
clarification. A lightweight XGBoost first stage handles clear cases, while
a selective LLM verifier is invoked only for uncertain, ambiguous, or
relation-heavy queries; a deterministic history resolver decides whether a
conversational reference can be resolved before retrieval. The system
operates over a dense index of legal text chunks and a large Vietnamese
legal knowledge graph built from the official electronic legal
codification (\textit{Bộ Pháp điển điện tử}). Through experiments on a
strict routing benchmark, an ambiguity diagnostic, and a conversation
stress test, we demonstrate (a)~that the two-stage router improves
end-to-end answer quality over an oracle with perfect routing labels,
indicating that the second stage acts as a Chain-of-Thought step;
(b)~that the first-stage classifier routes reliably at millisecond
latency and remains robust when surface templates are paraphrased away;
(c)~that clarification is detected with perfect precision and zero false
positives; and (d)~that selective second-stage invocation yields a
controllable quality--latency trade-off.

\smallskip
\noindent\textbf{Keywords:} Retrieval-Augmented Generation,
GraphRAG, Hybrid RAG, Adaptive Routing, Vietnamese Legal QA,
Knowledge Graphs, Multi-hop Reasoning, Ambiguity Detection.
\end{abstract}

% ============================================================
\section{Introduction}
\label{sec:intro}
% ============================================================

\subsection{Background and Motivation}

Retrieval-Augmented Generation (RAG) grounds language-model outputs in
external evidence, improving factual reliability without
retraining~\cite{lewis2020rag,guu2020realm}. Its most common form relies
on dense passage retrieval, embedding query and corpus into a shared
vector space and returning the most semantically similar
chunks~\cite{karpukhin2020dpr}. This is effective for factoid questions
answered in a single passage, but it treats the corpus as a flat
collection of independent chunks and reduces retrieval to nearest-neighbor
search over isolated spans~\cite{karpukhin2020dpr,trivedi2023interleaving}.

In practice, this assumption is overly simplistic for many real-world
question-answering settings. When the answer depends on relations among
entities, on document structure, or on chains of evidence spanning
multiple documents, dense retrieval cannot follow those connections
because the relations are not represented in the vector
space~\cite{trivedi2023interleaving}. Graph-based retrieval was introduced
precisely to address this gap: by organizing the corpus as a knowledge
graph, retrieval can traverse explicit, typed relationships rather than
relying solely on surface
similarity~\cite{edge2024graphrag,wang2024kgp,gutierrez2024hipporag}. Yet
graph traversal is more expensive than dense lookup, and applying it to
every query---including simple ones whose answer lies in a single
article---incurs unnecessary latency and computational cost. Vietnamese
legal question answering exposes this tension sharply: the same system
must serve both simple article-level definitions and multi-document
questions involving amendment, replacement, and cross-document
reference~\cite{chalkidis2020legalbert,zhong2020nlp_law,sovrano2020legalkg},
while Vietnamese-specific word segmentation and multi-syllable named
entities further complicate legal reference
matching~\cite{nguyen2020phobert,vu2018vncorenlp}.

When confronted with such heterogeneous queries, a naturally arising
question is how to decide \emph{which} retrieval strategy a given query
requires, \emph{before} any evidence is retrieved. Existing Hybrid RAG
systems sidestep this question by applying dense and graph retrieval
indiscriminately, paying the cost of both regardless of query
complexity~\cite{edge2024graphrag,he2024lightrag,gao2023modularrag}.
Recent adaptive-routing methods show that selecting a strategy by query
complexity improves both quality and
efficiency~\cite{jeong2024adaptive,chen2023frugalgpt,ong2024routellm}, but
they route among \emph{model sizes} or \emph{retrieval depths} rather than
among distinct \emph{evidence mechanisms}; the decision of whether a query
needs relational traversal at all is left unaddressed. A separate and
safety-critical gap concerns ambiguity. When a user asks
\textit{``Văn bản đó còn hiệu lực không?''} (\textit{``Is that document
still in effect?''}) without a usable referent in the conversation,
forcing retrieval grounds the answer in an arbitrary legal document and
returns a confidently wrong answer~\cite{min2020ambigqa,elgohary2019canard}.
No prior system simultaneously decides among dense retrieval, graph
traversal, hybrid reasoning, and clarification for Vietnamese legal QA.

Motivated by the observation that the appropriate retrieval strategy is a
property of the query itself---and can be predicted cheaply from the query
and its conversational context---we propose a two-stage Reasoning-aware
Adaptive Router that selects one of four evidence actions before
retrieval. A lightweight XGBoost classifier~\cite{chen2016xgboost} handles
clear cases at low cost; a selective LLM verifier is invoked only for
uncertain, ambiguous, or relation-heavy queries, following the cost-aware
cascade principle~\cite{chen2023frugalgpt,ong2024routellm}. A
deterministic \texttt{HistoryResolver} precedes both stages to distinguish
resolved, irrelevant, and conflicting conversation references, so that
clarification is triggered only when the legal target is genuinely
underspecified.

\subsection{Main Contributions}

To address the heterogeneity of legal queries described above, we adopt a
design principle in which the retrieval strategy is selected by a
two-stage router operating upstream of retrieval, rather than by running
all strategies and merging their outputs. The first stage performs a fast,
feature-based classification of the query into a retrieval mechanism using
an XGBoost model trained on routing labels derived from evidence
structure. Because this stage is lightweight, it can serve as a real-time
gateway that resolves the majority of clear cases without invoking a large
model. The second stage is an LLM verifier that is triggered selectively,
only when the first-stage decision is uncertain, when the query carries
strong relational signals, or when the conversational referent cannot be
resolved. This selective design keeps the expected per-query cost low
while reserving explicit reasoning for the cases that need it. Crucially,
the second stage does more than correct routing decisions: by reasoning
explicitly over query intent and legal constraints before answering, it
also functions as a Chain-of-Thought step that improves the quality of the
generated answer. Clarification is integrated as a first-class routing
action so that, when no retrieval strategy can safely ground the answer,
the system asks the user rather than guessing.

Our contributions are four-fold:
\begin{itemize}
  \item We formulate reasoning-aware retrieval routing for Vietnamese
  legal QA as a four-class decision---dense retrieval, graph traversal,
  hybrid reasoning, and clarification---made before evidence retrieval,
  unifying mechanism selection and ambiguity handling in a single
  framework~\cite{jeong2024adaptive,min2020ambigqa}.

  \item We design a two-stage router in which a lightweight XGBoost stage
  handles clear cases and a selective LLM verifier is invoked only for
  uncertain, ambiguous, or relation-heavy queries, and we show that this
  verifier acts as an explicit Chain-of-Thought mechanism: the Two-stage
  Hybrid (F1~$=0.661$) outperforms an Oracle Router with perfect routing
  labels (F1~$=0.606$), demonstrating that Stage~2 reasoning improves
  answer generation beyond correct route selection
  alone~\cite{chen2023frugalgpt}.

  \item We introduce a deterministic \texttt{HistoryResolver} that
  classifies a conversational referent as resolved, irrelevant, or
  conflicting, enabling the router to unblock retrieval when a referent is
  available and to force clarification when it is
  absent~\cite{elgohary2019canard}.

  \item Through experiments on a strict $600$-query benchmark, a
  $234$-query ambiguity benchmark, and a $160$-query conversation stress
  test, we validate (a)~end-to-end answer quality against fixed-strategy
  and oracle baselines, (b)~routing accuracy against rule-based and
  learned baselines, (c)~clarification behavior, and (d)~the cost-quality
  trade-off of selective verification.
\end{itemize}

To the best of our knowledge, this is the first system to treat
clarification as a first-class routing action alongside dense, graph, and
hybrid retrieval in a Hybrid GraphRAG pipeline for Vietnamese legal
question answering~\cite{sovrano2020legalkg,zhong2020nlp_law}.

\subsection{Organization}

Section~\ref{sec:related_work} reviews related work.
Section~\ref{sec:methodology} defines the routing problem, dataset,
and pipeline. Section~\ref{sec:routing_framework} details the routing
framework. Sections~\ref{sec:eval} and~\ref{sec:results} present
experimental setup and results. Section~\ref{sec:limitations}
discusses limitations before Section~\ref{sec:conclusion} concludes.

% ============================================================
\section{Related Work}
\label{sec:related_work}
% ============================================================

\textbf{Retrieval-Augmented Generation.}
RAG combines parametric knowledge from neural language models with
non-parametric knowledge retrieved from an external
corpus~\cite{lewis2020rag,guu2020realm}. Dense retrieval is effective for
open-domain factoid QA but degrades on multi-hop evidence
chains~\cite{karpukhin2020dpr,trivedi2023interleaving}. A line of advanced
variants improves retrieval control---Self-RAG through
self-reflection~\cite{asai2023selfrag}, FLARE through active
retrieval~\cite{jiang2023flare}, RAPTOR through hierarchical
summaries~\cite{sarthi2024raptor}, and CRAG through corrective
re-retrieval~\cite{yan2024corrective}. These methods refine \emph{how}
dense retrieval is performed, but they operate within a single evidence
mechanism and do not decide whether a query requires relational traversal
or clarification.

\textbf{Graph Retrieval and GraphRAG.}
Graph-based retrieval lets evidence follow explicit
relationships~\cite{sun2019pullnet,wang2024kgp}. GraphRAG builds entity
graphs for query-focused summarization~\cite{edge2024graphrag}; HippoRAG
introduces a neurobiologically inspired graph memory for associative
retrieval~\cite{gutierrez2024hipporag}; LightRAG uses dual-level retrieval
over entity and community knowledge~\cite{he2024lightrag}. For legal QA,
graph retrieval is especially relevant because legal documents encode
structured relations such as amendment, replacement, and
jurisdiction~\cite{sovrano2020legalkg}. However, these systems assume
graph traversal is always applied; none decides \emph{when} traversal is
warranted versus when a dense lookup would suffice at lower cost.

\textbf{Adaptive Routing and Cost-aware Inference.}
Adaptive-RAG selects a retrieval strategy according to question
complexity~\cite{jeong2024adaptive}, while FrugalGPT and RouteLLM route
queries across models of different
cost~\cite{chen2023frugalgpt,ong2024routellm}. These works establish the
cost-aware cascade principle that motivates our design. Their routing
target, however, is the \emph{model} or \emph{retrieval depth}; they do
not route among qualitatively different evidence mechanisms, and they do
not treat clarification as a routable action.

\textbf{Ambiguity Handling.}
CANARD studies question rewriting for context-dependent
questions~\cite{elgohary2019canard}, and AmbigQA shows that open-domain
questions often admit multiple valid
interpretations~\cite{min2020ambigqa}. These studies treat ambiguity
resolution as a stand-alone task applied before or after retrieval, rather
than as one possible outcome of a unified routing decision.

\textbf{Vietnamese NLP and Legal QA.}
VnCoreNLP provides Vietnamese word segmentation and
NER~\cite{vu2018vncorenlp}, and PhoBERT provides pretrained Vietnamese
representations~\cite{nguyen2020phobert}. Legal NLP more broadly requires
reasoning over regulatory
relations~\cite{chalkidis2020legalbert,zhong2020nlp_law,guha2023legalbench,louis2024interpretable}.
This body of work supplies the linguistic and legal-domain foundations our
system builds on, but does not address adaptive retrieval routing.

\paragraph{Discussion.}
Despite substantial progress along each of these lines, no prior system
decides among dense retrieval, graph traversal, hybrid reasoning, and
clarification \emph{before} retrieval in a single framework. Advanced-RAG
methods improve a single mechanism; GraphRAG methods assume traversal is
always needed; cost-aware routers select models rather than evidence
mechanisms; and ambiguity-handling methods treat clarification in
isolation. Our work differs in placing an explicit, reasoning-aware
routing decision at the center of the pipeline, and in treating
clarification as a first-class route backed by deterministic
conversational context resolution.
\FloatBarrier

% ============================================================
\section{Methodology}
\label{sec:methodology}
% ============================================================

\subsection{Problem Definition}
\label{subsec:problem_definition}

Let $\mathcal{C}=\{D_1,\ldots,D_N\}$ denote the Vietnamese legal
corpus. From $\mathcal{C}$, the system builds two complementary
knowledge sources: a dense vector index
$\mathcal{I}_v=\{(c_j,\mathbf{e}_j)\}_{j=1}^{M}$ of $M$ legal text
chunks, and a legal knowledge graph
$\mathcal{G}=(\mathcal{V},\mathcal{E})$ where a node
$v\in\mathcal{V}$ represents a legal document, article, semantic concept,
agency, or other entity, and an edge $(u,r,v)\in\mathcal{E}$ represents a
typed relation such as \textsc{HAS\_ARTICLE}, \textsc{MENTIONS},
\textsc{AMENDS}, or \textsc{REPLACES}.

At inference time, the input is the current query $q_t$ at turn
$t\geq 1$, together with optional conversation history
$H=\{(q_i,a_i)\}_{i=1}^{t-1}$. When $t=1$, $H=\emptyset$. The
router selects one action from the route set
$\mathcal{D}=\{d_v,\,d_g,\,d_h,\,d_c\}$, where $d_v$, $d_g$,
$d_h$, $d_c$ denote \texttt{dense\_retrieval},
\texttt{graph\_traversal}, \texttt{hybrid\_reasoning}, and
\texttt{clarify}, respectively. The routing function is
\begin{equation}
\rho:(q_t,H)\rightarrow d^*,\qquad d^*\in\mathcal{D}.
\label{eq:routing_function}
\end{equation}
For retrieval routes, the evidence set is
\begin{equation}
E_{d^*}=\begin{cases}
\mathrm{TopK}(q_t,\mathcal{I}_v,k), & d^*=d_v,\\
\mathrm{Traverse}(q_t,\mathcal{G},h), & d^*=d_g,\\
E_v\cup E_g, & d^*=d_h,
\end{cases}
\label{eq:evidence_set}
\end{equation}
where $k$ is the dense retrieval depth and $h$ is the traversal
depth. The generator produces $a_t=\mathrm{Generate}(q_t,H,E_{d^*})$.
If $d^*=d_c$, no evidence is retrieved; instead the system returns
a clarification question $c_t=\mathrm{Clarify}(q_t,H)$.

\subsection{Dataset}
\label{subsec:dataset}

The evaluation uses a self-constructed Vietnamese legal QA routing
dataset derived from a single source: the Vietnamese electronic legal
codification system (\textit{Bộ Pháp điển điện tử}), flattened into
article-level records. QA pairs are linked to supporting articles through
an \texttt{article\_key} (format \texttt{DOC::ARTICLE\_ID}) and assigned
routing labels based on evidence structure: \texttt{dense\_retrieval} for
single-article questions, \texttt{graph\_traversal} for multi-article
reasoning within one document, and \texttt{hybrid\_reasoning} for evidence
spanning multiple documents. Low-quality samples are removed by
deduplication, length filtering, and placeholder-answer filtering. To
reduce template-induced lexical leakage in the routing features, the
surface forms of the generated queries are additionally paraphrased while
preserving their underlying evidence structure, so that the routing
classifier must rely on genuine structural signals rather than on
memorized template wording.

Table~\ref{tab:dataset} summarizes the dataset splits. The dataset
remains skewed toward dense retrieval, which accounts for half of the
samples ($50\%$), motivating Macro-F1 reporting alongside accuracy. The
strict split is used for retrieval-route evaluation; the \texttt{clarify}
route is evaluated separately on a $234$-query ambiguity benchmark
($156$ ambiguous, $78$ answerable controls) and a $160$-query conversation
stress test.

\begin{table}[!htbp]
\centering
\caption{Vietnamese legal QA routing dataset split.}
\label{tab:dataset}
\small
\begin{tabular}{lrrrr}
\toprule
\textbf{Split} & \textbf{Dense} & \textbf{Graph} &
\textbf{Hybrid} & \textbf{Total} \\
\midrule
Train & 400 & 200 & 200 & 800 \\
Dev & 50 & 25 & 25 & 100 \\
Test & 300 & 150 & 150 & 600 \\
\midrule
\textbf{Total} & 750 & 375 & 375 & 1,500 \\
\bottomrule
\end{tabular}
\end{table}

\subsection{System Overview}
\label{subsec:system_overview}

\begin{figure}[!htbp]
\centering
\includegraphics[width=0.97\columnwidth]{figs/system_architecture}
\caption{High-level data flow. The Vietnamese legal corpus
$\mathcal{C}$ is processed offline into a dense vector index
$\mathcal{I}_v$ (ChromaDB) and a legal knowledge graph $\mathcal{G}$
(Neo4j). At query time, the two-stage adaptive router selects evidence
from $\mathcal{I}_v$, $\mathcal{G}$, or both, or returns a
clarification question $c_t$.}
\label{fig:system_architecture}
\end{figure}

\begin{figure*}[!htbp]
\centering
\includegraphics[width=\textwidth]{figs/system_pipeline}
\caption{Reasoning-aware adaptive routing pipeline and illustrative
routing outcomes. (a)~The pipeline processes $(q_t,H)$ through feature
extraction, Stage~1 routing, trigger checking, optional Stage~2 LLM
verification, and route-controlled retrieval or clarification.
(b)~Concrete examples: dense ($d_v$) for direct lookup, graph ($d_g$)
for amendment relations, hybrid ($d_h$) for cross-document synthesis,
and clarify ($d_c$) for underspecified referents.}
\label{fig:system_pipeline}
\end{figure*}

Figure~\ref{fig:system_architecture} shows the overall data flow. In
the offline phase, $\mathcal{C}$ is indexed once into $\mathcal{I}_v$
and $\mathcal{G}$. In the online phase, every $(q_t,H)$ is routed
against these fixed sources. Figure~\ref{fig:system_pipeline}
illustrates both the routing pipeline and representative examples of
the four routing actions. The final route $d^*$ controls evidence
retrieval, graph reasoning, hybrid evidence aggregation, or
clarification generation.

\subsection{Implementation Details}
\label{subsec:implementation}

The Chroma vector store is persisted with chunks embedded by
\path{microsoft/Harrier-OSS-v1-0.6B} (256-token chunks, 32-token
overlap, top-$k{=}8$). The graph backend uses Neo4j with full-text
indexes over legal documents, articles, and vector chunks; graph
retrieval uses Text-to-Cypher querying via an OpenAI-compatible
endpoint, generating Cypher queries against $\mathcal{G}$. Stage~2
uses a local vLLM endpoint serving a Qwen3-32B-AWQ model, prompted
for structured JSON route prediction (route decision, confidence,
optional clarification question, and reasoning trace). Vietnamese NER
uses \path{NlpHUST/ner-vietnamese-electra-base}~\cite{clark2020electra}.
The Stage~1 router checkpoint is stored at
\path{data/router\_training/legal\_strict/router\_model.pkl}.

\FloatBarrier

% ============================================================
\section{Reasoning-aware Adaptive Routing Framework}
\label{sec:routing_framework}
% ============================================================

\subsection{Legal Knowledge Base}
\label{subsec:legal_kb}

The vector index $\mathcal{I}_v$ contains $68{,}835$ embedded chunks
appropriate for direct questions whose answers are stated in legal text.
The graph $\mathcal{G}$ stores explicit typed relations among legal
documents, articles, and concepts.
Tables~\ref{tab:graph_stats} and~\ref{tab:graph_relations} report the
statistics of the graph built from the Pháp~điển codification, in which
article nodes are promoted to \textsc{LegalArticle} to enable
article-level evidence retrieval.

\begin{table}[!htbp]
\centering
\caption{Vietnamese legal graph $\mathcal{G}$ statistics.}
\label{tab:graph_stats}
\small
\begin{tabular}{lr}
\toprule
\textbf{Graph item} & \textbf{Count} \\
\midrule
Nodes $|\mathcal{V}|$ & 804,322 \\
Edges $|\mathcal{E}|$ & 1,968,621 \\
Legal documents & 4,598 \\
Legal articles (\textsc{LegalArticle}) & 68,570 \\
Semantic concepts & 392,592 \\
Other entities & 338,562 \\
\bottomrule
\end{tabular}
\end{table}

\begin{table}[!htbp]
\centering
\caption{Main relation types in the legal graph (top by count).}
\label{tab:graph_relations}
\small
\begin{tabular}{lr}
\toprule
\textbf{Relation type} & \textbf{Count} \\
\midrule
\textsc{MENTIONS} & 1,033,341 \\
\textsc{HAS\_CONDITION} & 202,298 \\
\textsc{OBLIGATES} & 128,768 \\
\textsc{PART\_OF\_PROCEDURE} & 102,004 \\
\textsc{REGULATES} & 75,137 \\
\bottomrule
\end{tabular}
\end{table}

The graph is large and densely linked at the semantic level through
relations such as \textsc{MENTIONS} ($1{,}033{,}341$ edges) and
\textsc{HAS\_CONDITION} ($202{,}298$), whereas explicit cross-document
legal-effect relations such as \textsc{AMENDS} ($5{,}944$ edges) and
\textsc{REPLACES} ($4{,}189$) remain comparatively sparse. This asymmetry
directly motivates adaptive routing: graph traversal is valuable for
relation-heavy queries but unnecessarily costly for simple article-level
lookups.

\subsection{Reasoning-aware Feature Extraction}
\label{subsec:features}

For each input $(q,H)$, the router extracts a feature vector
$\mathbf{f}(q,H)$ of $27$ features that blend lexical signals with
reasoning-aware metrics, together with two auxiliary scores: the ambiguity
score $\phi_{\mathrm{amb}}(q,H)\in[0,1]$ and the multi-hop score
$\xi(q)\in\{0,1,2\}$. The Stage~1 classifier relies most heavily on a
small set of structural features. The top components driving the routing
decision are: (1)~the \textbf{legal reference count} (feature importance
$0.3225$), which tracks how many distinct legal documents are mentioned;
(2)~the \textbf{query length in words} ($0.2131$), as cross-document
synthesis questions tend to be longer; (3)~\textbf{cross-document signals}
($0.0945$), reflecting explicit synthesis indicators; and (4)~the
\textbf{article reference count} ($0.0917$). Additional reasoning metrics
such as multi-hop connectors, sub-question count, procedural depth, and a
discrete complexity level provide finer-grained separation for harder
cases. Question-type cues (\textit{gì}, \textit{khi nào},
\textit{bao nhiêu}) and legal-relation cues (\textit{sửa đổi},
\textit{thay thế}, \textit{có hiệu lực từ}) further inform the decision.

The multi-hop score $\xi(q)=0$ means no relation signal is detected
(likely a direct lookup); $\xi(q)=1$ means one relation cue is present;
$\xi(q)=2$ means two or more signals co-occur, indicating the strongest
multi-hop demand. A deterministic \texttt{HistoryResolver} runs before
ambiguity scoring:
\begin{equation}
r_H = \mathrm{Resolve}(q_t, H),
\label{eq:history_resolver}
\end{equation}
returning candidate referents and a
\texttt{history\_resolution\_status}
$\in\{$\texttt{resolved}, \texttt{no\_history},
\texttt{irrelevant\_history}, \texttt{conflicting\_history},
\texttt{not\_needed}$\}$. If \textit{văn bản đó} has a unique
document referent in history, the resolved referent can unblock graph
traversal; otherwise the same surface form remains ambiguous.

\subsection{Stage~1 XGBoost Router}
\label{subsec:stage1}

Stage~1 is an XGBoost multi-class classifier~\cite{chen2016xgboost}
trained on the strict three-label routing data:
\begin{equation}
(d_1, s_1) = f_1\bigl(\mathbf{f}(q,H)\bigr),
\label{eq:stage1}
\end{equation}
where $d_1\in\{d_v,d_g,d_h\}$ is the initial route and
$s_1\in[0,1]$ is the confidence. Because the training split contains
only \texttt{dense}, \texttt{graph}, and \texttt{hybrid} labels,
Stage~1 cannot directly learn \texttt{clarify}---an architectural
property, not a bug. XGBoost is selected for fast, CPU-friendly
routing with reliable confidence scores. It runs at
$\approx\!8.9$~ms per query, versus $\approx\!53.3$~ms for a
PhoBERT~\cite{nguyen2020phobert} classifier, making it suitable as a
real-time gateway; this motivates a future three-level cascade
XGBoost~$\rightarrow$~PhoBERT~$\rightarrow$~LLM following
FrugalGPT~\cite{chen2023frugalgpt}.

\subsection{Stage~2 LLM Verifier}
\label{subsec:stage2}

Stage~2 is an LLM verifier $f_2$ invoked selectively. Three binary
trigger indicators are defined:
\begin{align}
\gamma_{\mathrm{conf}} &= \mathbb{I}[s_1 < \tau_{\mathrm{conf}}],
    \label{eq:gamma_conf}\\
\gamma_{\mathrm{amb}} &=
    \mathbb{I}[\phi_{\mathrm{amb}}(q,H) > \tau_{\mathrm{force}}],
    \label{eq:gamma_amb}\\
\gamma_{\mathrm{hop}} &=
    \mathbb{I}[d_1\neq d_g \wedge \xi(q)=2].
    \label{eq:gamma_hop}
\end{align}
The composite trigger is
$\Gamma(q,H,r_H) = \mathbb{I}[\gamma_{\mathrm{conf}} \vee
\gamma_{\mathrm{amb}} \vee \gamma_{\mathrm{hop}}]$,
and the final route is
\begin{equation}
d^* = \begin{cases}
f_2(q,H,d_1), & \Gamma(q,H,r_H)=1,\\
d_1, & \Gamma(q,H,r_H)=0.
\end{cases}
\label{eq:final_route}
\end{equation}
Default thresholds are $\tau_{\mathrm{conf}}=0.50$,
$\tau_{\mathrm{force}}=0.60$, and clarification override
$\tau_{\mathrm{clar}}=0.80$. In words, Stage~2 is invoked on low-confidence
Stage~1 decisions, on high ambiguity risk, or on strong multi-hop signals;
a \texttt{resolved} history unblocks retrieval and suppresses false
clarification, whereas missing, irrelevant, or conflicting history forces
verification or clarification.

Algorithm~\ref{alg:router} presents the routing procedure. The
deployed implementation additionally uses heuristics: a dense fast
path that skips Stage~2 for high-confidence low-ambiguity dense
queries, and a reasoning-signal check that escalates queries with
strong legal features even at moderate Stage~1 confidence.

\begin{algorithm}[!htbp]
\caption{Two-stage Adaptive Router}
\label{alg:router}
\begin{algorithmic}[1]
\Require $q$, $H$,
    $\tau_{\mathrm{conf}},\tau_{\mathrm{force}},\tau_{\mathrm{clar}}$
\Ensure Final route $d^*\in\mathcal{D}$
\State $r_H\leftarrow\mathrm{Resolve}(q,H)$
\State Compute $\mathbf{f}(q,H,r_H)$,
       $\phi_{\mathrm{amb}}(q,H,r_H)$, $\xi(q)$
\State $(d_1,s_1)\leftarrow f_1(\mathbf{f}(q,H,r_H))$
\State Compute $\Gamma(q,H,r_H)$
\If{$\Gamma(q,H,r_H)=1$}
  \State $d^*\leftarrow f_2(q,H,d_1)$
\Else
  \State $d^*\leftarrow d_1$
\EndIf
\State \Return $d^*$
\end{algorithmic}
\end{algorithm}

\paragraph{Running Example.}
For $q_1=$ ``\textit{Điều kiện để được cấp giấy phép xây dựng là
gì?}'', $\xi(q_1)=0$ and Stage~1 predicts $d_v$ with high confidence;
Stage~2 is bypassed. For $q_2=$ ``\textit{Nghị định nào sửa đổi quy
định về xử phạt vi phạm hành chính trong lĩnh vực đất đai?}'', the
phrase \textit{sửa đổi} raises legal-relation signals; the router
selects $d_g$ or $d_h$. For $q_3=$ ``\textit{Văn bản đó còn hiệu lực
không?}'' with $H=\emptyset$, the referent is unresolved,
$\phi_{\mathrm{amb}}$ is high, $\Gamma=1$, and $f_2$ returns
$d^*=d_c$: ``\textit{Bạn muốn hỏi hiệu lực của văn bản nào?}''.

\subsection{Route Execution and Fallback Policy}
\label{subsec:route_execution}

For $d_v$, the system retrieves top-$k$ chunks from $\mathcal{I}_v$.
For $d_g$, it links query entities to graph nodes and traverses
$\mathcal{G}$. For $d_h$, it constructs $E_{d^*}=E_v\cup E_g$. For
$d_c$, it returns a clarification question without retrieval. A
defensive fallback policy applies when primary retrieval is
uninformative: a failed graph traversal or hybrid synthesis falls back
to dense retrieval, an uninformative dense lookup may retry on the
graph, and a clarify route falls back to the ambiguity-detector's
default question.

\FloatBarrier

% ============================================================
\section{Experimental Evaluation}
\label{sec:eval}
% ============================================================

\subsection{Research Questions}

\textbf{RQ1} (\textit{End-to-end quality}): Does adaptive routing
improve answer quality over fixed-strategy and oracle baselines? Does
Stage~2 contribute beyond mere routing correction?
\textbf{RQ2} (\textit{Routing accuracy}): How accurately does the
Stage~1 classifier select routes, compared with alternative routing
approaches?
\textbf{RQ3} (\textit{Ambiguity handling}): Can the system avoid
unsafe retrieval by clarifying when the legal target is underspecified?
\textbf{RQ4} (\textit{Threshold ablation}): How do the key routing
thresholds affect routing accuracy and Stage~2 trigger rate?
\textbf{RQ5} (\textit{Cost-quality}): Does selective Stage~2
invocation provide a better trade-off than always-on or always-off
verification?

\subsection{Systems and Baselines}

Table~\ref{tab:systems} defines all compared systems. The
\textbf{Oracle Router} uses gold route labels to dispatch each query
to the gold-labeled retrieval path with identical retrieval depth and
generation settings as all other end-to-end systems. Comparing the
oracle against the Two-stage Hybrid directly quantifies how much answer
quality is attributable to Stage~2 reasoning beyond correct routing.

\begin{table}[!htbp]
\centering
\caption{Systems compared in this work.}
\label{tab:systems}
\scriptsize
\setlength{\tabcolsep}{3pt}
\renewcommand{\arraystretch}{1.08}
\begin{tabularx}{\columnwidth}{L{2.0cm}Y}
\toprule
\textbf{System} & \textbf{Description} \\
\midrule
Pure Vector & Always dense retrieval from Chroma. \\
Pure Graph & Always graph traversal in Neo4j (Text-to-Cypher). \\
Pure Hybrid & Always combines dense + graph evidence. \\
Single-stage Router & XGBoost Stage~1 only; no LLM verifier. \\
Two-stage Hybrid & Stage~1 + selective Stage~2 (full system). \\
Oracle Router & Gold route labels; same retrieval and generation. \\
\bottomrule
\end{tabularx}
\end{table}

\subsection{Evaluation Metrics}
\label{subsec:metrics}

\textbf{Answer quality:} Exact Match (EM) and token-level F1 computed
by \texttt{evaluate\_prediction(ground\_truth, answer)}, identical
across all end-to-end systems. Hit@$k$ measures retrieval coverage
(retrieved source matches gold within top-$k$).

\textbf{Route classification:} Accuracy, Macro-F1, Weighted-F1, and
per-class F1. Macro-F1 is reported because the dataset is dominated by
the dense class~\cite{powers2011evaluation}.

\textbf{Clarification:} Precision, recall, and F1 for \texttt{clarify}
on the ambiguity benchmarks.

\textbf{Cost-quality:} Average latency, p95, Stage~2 trigger rate,
and graph-to-dense fallback rate.

% ============================================================
\section{Results and Discussion}
\label{sec:results}
% ============================================================

\subsection{RQ1: End-to-end Answer Quality}
\label{subsec:rq1}

Since the ultimate purpose of routing is to produce better answers, we
first ask whether adaptive routing improves end-to-end answer quality over
fixed-strategy baselines, and whether the Stage~2 verifier contributes
beyond merely correcting route selection. To isolate the latter effect, we
compare against an Oracle Router that dispatches every query using its gold
route label. Table~\ref{tab:end_to_end} reports results on the strict
$600$-query benchmark, and Figure~\ref{fig:f1_bar} visualizes Answer F1
across systems.

\begin{table*}[!htbp]
\centering
\caption{End-to-end results on the strict $600$-query Vietnamese legal
QA benchmark. F1 and EM are computed by
\texttt{evaluate\_prediction(gold\_context,~answer)}, identical across
all systems. Hit@$k$ measures retrieval coverage on $N{=}519$ queries
(excluding \textit{Toàn văn} queries). S2~=~Stage~2 trigger rate;
FB~=~graph-to-dense fallback rate. Bold marks best
per column, excluding Oracle.}
\label{tab:end_to_end}
\scriptsize
\setlength{\tabcolsep}{2.5pt}
\renewcommand{\arraystretch}{1.05}
\begin{tabular}{lrrrrrrrr}
\toprule
\textbf{System} & \textbf{F1} & \textbf{EM} & \textbf{Hit@1} &
\textbf{Hit@3} & \textbf{Lat (ms)} & \textbf{p95 (ms)} &
\textbf{S2} & \textbf{FB} \\
\midrule
Pure Vector  & \textbf{0.701} & 0.170 & \textbf{0.588} & \textbf{0.697} &
               5,682 & 16,895 & -- & -- \\
Pure Graph   & 0.632 & \textbf{0.327} & 0.542 & 0.653 &
               5,588 & 17,332 & -- & -- \\
Pure Hybrid  & 0.656 & 0.098 & 0.003 & 0.007 &
               6,815 & 20,174 & -- & -- \\
Single-stage & 0.530 & 0.140 & 0.292 & 0.362 &
               \textbf{4,749} & \textbf{15,953} & -- & 3.7\% \\
Two-stage (ours) & 0.661 & 0.267 & 0.410 & 0.473 &
               9,950 & 23,686 & 78.3\% & 5.3\% \\
\midrule
Oracle       & 0.606 & 0.127 & 0.282 & 0.360 &
               5,975 & 19,330 & -- & -- \\
\bottomrule
\end{tabular}
\end{table*}

\begin{figure}[!htbp]
\centering
\includegraphics[width=1.0\columnwidth]{figs/f1_bar}
\caption{Answer F1 on the strict $600$-query benchmark. Pure Vector is
the strongest fixed-strategy baseline ($0.701$). The Two-stage Router
($0.661$) outperforms the Oracle ($0.606$), confirming Stage~2
reasoning adds genuine value beyond routing-label quality. The
Single-stage Router ($0.530$) is the weakest, underscoring the
critical role of Stage~2.}
\label{fig:f1_bar}
\end{figure}

\paragraph{Interpretation.}
We make the following observations:
\begin{itemize}
  \item \textbf{Stage~2 acts as a Chain-of-Thought mechanism, not merely
  a route corrector.} The Two-stage Hybrid (F1~$=0.661$) beats the
  Oracle Router (F1~$=0.606$) by $+0.055$ despite the Oracle's gold labels.
  Since both dispatch correct queries to identical paths, the gain comes
  from Stage~2's explicit reasoning over query intent and legal
  constraints, improving generation beyond route selection alone.

  \item \textbf{Stage~2 accounts for the entire routing-to-quality gain.}
  The gap between the Single-stage Router (F1~$=0.530$) and the Two-stage
  Hybrid (F1~$=0.661$) is $+0.131$, attributable entirely to selective
  Stage~2 verification, since the two systems share an identical Stage~1.

  \item \textbf{Pure Vector is the strongest fixed-strategy baseline.}
  It attains the highest fixed-strategy F1 ($0.701$) and coverage
  (Hit@1~$=0.588$): the dense index covers direct-lookup queries well. The
  router's value lies not in beating dense retrieval everywhere but in the
  queries it cannot answer---and in surpassing the oracle via Stage~2.

  \item \textbf{Pure Graph trades fluency for relational precision.} It
  reaches the highest EM ($0.327$) on relation-heavy queries but lower
  F1 ($0.632$), reflecting the less fluent answers produced from
  structured graph context.

  \item \textbf{Pure Hybrid exposes a concrete engineering gap.} Despite
  a competitive F1 ($0.656$), its retrieval coverage collapses to near
  zero (Hit@1~$=0.003$) due to a source-identifier format mismatch between
  the dense and graph indices---a gap we return to in
  Section~\ref{sec:limitations}.
\end{itemize}

\subsection{RQ2: Routing Accuracy}
\label{subsec:rq2}

A router is only useful if it selects the correct retrieval mechanism. We
therefore evaluate the offline routing accuracy of the Stage~1 classifier
on the paraphrased strict test split and compare it against rule-based and
learned baselines (Table~\ref{tab:routing_baselines}).

\begin{table}[!htbp]
\centering
\caption{RQ2: Routing classifier comparison on the paraphrased strict
test split.}
\label{tab:routing_baselines}
\scriptsize
\setlength{\tabcolsep}{3pt}
\renewcommand{\arraystretch}{1.08}
\begin{tabular}{lrrr}
\toprule
\textbf{Classifier} & \textbf{Acc.} & \textbf{Mac-F1} & \textbf{Wt-F1} \\
\midrule
MajorityRoute & 0.500 & 0.222 & 0.333 \\
KeywordRuleRouter & 0.528 & 0.298 & 0.398 \\
Simple 3-Rule Router & 0.780 & 0.670 & -- \\
Logistic Regression & 0.985 & 0.982 & 0.985 \\
Random Forest & 0.992 & 0.990 & 0.992 \\
XGBoost (Stage~1) & \textbf{0.995} & \textbf{0.995} & \textbf{0.995} \\
\bottomrule
\end{tabular}
\end{table}

We make the following observations:
\begin{itemize}
  \item \textbf{Learned classifiers dominate heuristic baselines.} The
  majority-class router (Macro-F1~$=0.222$), the keyword rule router
  ($0.298$), and even a hand-tuned three-rule router ($0.670$) fall well
  below the learned models, confirming that routing among evidence
  mechanisms cannot be reduced to a handful of surface rules.

  \item \textbf{XGBoost remains accurate after paraphrasing.} XGBoost
  reaches Macro-F1~$=0.995$ on the test split whose surface forms have
  been paraphrased to remove template wording. Because the rigid template
  cues that would otherwise leak the label are no longer present, this
  high score indicates that the classifier relies on genuine structural
  features rather than memorized templates.

  \item \textbf{Feature importance confirms reliance on structural
  signals.} The decision is driven mainly by the legal reference count
  ($0.3225$), query length ($0.2131$), cross-document signals ($0.0945$),
  and article reference count ($0.0917$)---all meaningful indicators of
  evidence structure rather than spurious lexical artifacts.

  \item \textbf{XGBoost is an efficient real-time gateway.} It runs in
  $\approx\!8.9$~ms on CPU versus $\approx\!53.3$~ms on GPU for a PhoBERT
  classifier, making it well suited to gating the expensive Stage~2
  verifier. The decisive evidence for the system, however, is the
  end-to-end answer-quality and latency trade-off of RQ1 and RQ5, where
  the full pipeline must execute the selected reasoning paths.
\end{itemize}

\subsection{RQ3: Ambiguity and Clarification Handling}
\label{subsec:rq3}

Routing is safety-critical when the legal target is underspecified:
retrieving against an arbitrary document yields a confidently wrong
answer. We therefore evaluate whether the system clarifies instead of
guessing, on the $234$-query ambiguity benchmark
(Table~\ref{tab:ambiguity}) and broken down by ambiguity type
(Table~\ref{tab:ambiguity_type}).

\begin{table}[!htbp]
\centering
\caption{RQ3: Clarification on the $234$-query ambiguity benchmark.}
\label{tab:ambiguity}
\scriptsize
\setlength{\tabcolsep}{2.5pt}
\resizebox{\columnwidth}{!}{%
\begin{tabular}{lrrrrr}
\toprule
\textbf{Variant} & \textbf{Route Acc.} & \textbf{P} &
\textbf{R} & \textbf{F1} & \textbf{S2 Trig.} \\
\midrule
Stage~1 only & 0.269 & 0.000 & 0.000 & 0.000 & 0.000 \\
Stage~1 + Stage~2 (initial) & 0.585 & 1.000 & 0.519 & 0.684 & 0.577 \\
Phase~3 full policy & \textbf{0.722} & \textbf{1.000} &
\textbf{0.724} & \textbf{0.840} & -- \\
\bottomrule
\end{tabular}%
}
\end{table}

\begin{table}[!htbp]
\centering
\caption{RQ3: Stage~2 trigger rate and recall by ambiguity type
(initial Stage~1+Stage~2 configuration).}
\label{tab:ambiguity_type}
\small
\setlength{\tabcolsep}{3.0pt}
\begin{tabular}{lrrr}
\toprule
\textbf{Ambiguity type} & \textbf{Total} & \textbf{S2 Trig.} & \textbf{Recall} \\
\midrule
Pronoun reference & 42 & 1.000 & 1.000 \\
Incomplete context & 39 & 1.000 & 1.000 \\
Missing entity & 39 & 0.000 & 0.000 \\
Multi-interpretation & 36 & 0.000 & 0.000 \\
\bottomrule
\end{tabular}
\end{table}

We make the following observations:
\begin{itemize}
  \item \textbf{Stage~1 alone cannot clarify.} It scores clarify
  F1~$=0.000$ because it is trained on only three retrieval labels; the
  clarification capability is therefore an emergent property of the
  Stage~2 verifier and ambiguity detector, not of the base classifier.

  \item \textbf{Stage~2 clarifies with perfect precision.} The initial
  Stage~1+Stage~2 configuration raises clarify F1 to $0.684$ at precision
  $1.000$: every query the system chooses to clarify is genuinely
  ambiguous, so clarification never disrupts an answerable query.

  \item \textbf{History resolution and a severe-ambiguity override close
  most of the recall gap.} Adding the Phase~3 \texttt{HistoryResolver} and
  override raises clarify F1 to $0.840$ while preserving zero false
  positives, improving recall from $0.519$ to $0.724$.

  \item \textbf{Surface-contextual ambiguity is solved; semantic
  ambiguity is not.} Pronoun-reference and incomplete-context queries are
  detected with trigger rate and recall both $1.000$, whereas
  missing-entity and multi-interpretation queries are missed entirely
  ($0.000$): the former look grammatically complete, and the latter admit
  several valid legal readings. These two categories are the principal
  targets for future work.

  \item \textbf{The benefit transfers to multi-turn conversation.} On the
  separate $160$-query conversation stress test, clarify F1 improves from
  $0.545$ before Phase~3 to $0.862$ after, with history-resolution
  accuracy $0.850$. For instance, given history mentioning
  ``\textit{Nghị định 100/2019/NĐ-CP}'', the follow-up
  ``\textit{Văn bản đó còn hiệu lực không?}'' is resolved and routed to
  graph traversal, whereas with empty history the same query is routed to
  \texttt{clarify}.
\end{itemize}

\subsection{RQ4: Threshold Ablation}
\label{subsec:rq4}

The router exposes two key thresholds: the confidence threshold
$\tau_{\mathrm{conf}}$ that gates Stage~2, and the clarification override
threshold $\tau_{\mathrm{clar}}$. We study their effect on the
development set in Table~\ref{tab:threshold_ablation}.

\begin{table}[!htbp]
\centering
\caption{RQ4: Threshold ablation on the development set.}
\label{tab:threshold_ablation}
\small
\setlength{\tabcolsep}{3.5pt}
\begin{tabular}{lrrr}
\toprule
\textbf{Parameter} & \textbf{Value} & \textbf{Acc.} & \textbf{S2 rate} \\
\midrule
$\tau_{\mathrm{conf}}$ & 0.50 & 0.7529 & 0.2299 \\
$\tau_{\mathrm{conf}}$ & 0.70 & 0.7529 & 0.3046 \\
$\tau_{\mathrm{conf}}$ & 0.95 & 0.7529 & 0.5920 \\
$\tau_{\mathrm{conf}}$ & 0.99 & 0.7529 & 0.5977 \\
\midrule
$\tau_{\mathrm{clar}}$ & 0.50 & 0.6839 & 0.2299 \\
$\tau_{\mathrm{clar}}$ & 0.70 & 0.7529 & 0.2299 \\
$\tau_{\mathrm{clar}}$ & 0.90 & 0.7759 & 0.2299 \\
$\tau_{\mathrm{clar}}$ & 0.95 & 0.8448 & 0.2299 \\
\bottomrule
\end{tabular}
\end{table}

We make the following observations:
\begin{itemize}
  \item \textbf{$\tau_{\mathrm{conf}}$ governs cost, not accuracy.}
  Sweeping $\tau_{\mathrm{conf}}$ from $0.50$ to $0.99$ leaves routing
  accuracy unchanged at $0.7529$ while raising the Stage~2 trigger rate
  from $0.2299$ to $0.5977$. Many Stage~1 predictions are already
  high-confidence, so $\tau_{\mathrm{conf}}$ controls how much
  verification is spent rather than how accurate the routing is.

  \item \textbf{$\tau_{\mathrm{clar}}$ has a strong accuracy effect.} A
  low value of $0.50$ over-clarifies on the non-ambiguous development set
  (accuracy $0.6839$), whereas $0.95$ raises accuracy to $0.8448$, since a
  stricter override prevents spurious clarification of answerable queries.

  \item \textbf{The two benchmarks pull in opposite directions.} This
  ablation uses routing-oriented development data, not the ambiguity
  benchmark; the $\tau_{\mathrm{clar}}$ that maximizes strict routing
  accuracy is not the one that maximizes clarify recall, so the two
  objectives must be jointly calibrated rather than tuned in isolation.
\end{itemize}

\subsection{RQ5: Cost-Quality Trade-off}
\label{subsec:rq5}

Stage~2 is powerful but expensive, so the central design question is
whether selective invocation pays off relative to always-off and always-on
verification. Table~\ref{tab:cost} reports the latency-quality analysis.

\begin{table}[!htbp]
\centering
\caption{RQ5: Cost-quality trade-off. S2~=~Stage~2 trigger rate.
Oracle latency reflects gold-label routing without Stage~2 overhead.}
\label{tab:cost}
\scriptsize
\setlength{\tabcolsep}{3pt}
\renewcommand{\arraystretch}{1.05}
\begin{tabular}{lrrrrr}
\toprule
\textbf{System} & \textbf{F1} & \textbf{Avg (ms)} &
\textbf{p95 (ms)} & \textbf{S2} & \textbf{FB} \\
\midrule
Pure Vector & 0.701 & 5,682 & 16,895 & 0\% & 0\% \\
Pure Graph & 0.632 & 5,588 & 17,332 & 0\% & 0\% \\
Pure Hybrid & 0.656 & 6,815 & 20,174 & 0\% & 0\% \\
Single-stage & 0.530 & \textbf{4,749} & \textbf{15,953} & 0\% & 3.7\% \\
Two-stage & \textbf{0.661} & 9,950 & 23,686 & 78.3\% & 5.3\% \\
\midrule
Oracle & 0.606 & 5,975 & 19,330 & 0\% & 0\% \\
\bottomrule
\end{tabular}
\end{table}

We make the following observations:
\begin{itemize}
  \item \textbf{Always-off verification is cheap but weak.} The
  Single-stage Router is fastest ($4{,}749$~ms) by skipping Stage~2
  entirely, but pays the largest quality penalty (F1~$=0.530$).

  \item \textbf{Stage~2 is two orders of magnitude more expensive than
  Stage~1.} Routing-only measurements show Stage~2 calls cost
  $\approx\!4{,}017$~ms versus $\approx\!33.7$~ms for Stage~1-only
  routing---a roughly $119\times$ increase---so always-on verification is
  not viable; selective invocation is the only practical design. At the
  observed $78.3\%$ trigger rate, the Two-stage system runs at
  $9{,}950$~ms average.

  \item \textbf{The overhead comes from verification, not routing, and is
  justified.} The Oracle Router runs at $5{,}975$~ms with no Stage~2
  overhead; the Two-stage Hybrid adds about $3{,}975$~ms on top, which is
  attributable to selective Stage~2 invocation rather than to routing
  itself, and is repaid by the $+0.055$~F1 gain over the Oracle. Lowering
  $\tau_{\mathrm{conf}}$ (RQ4) reduces the trigger rate and hence this
  overhead, giving a tunable quality--latency dial.
\end{itemize}

\subsection{Overall Discussion}
\label{subsec:discussion}

The experiments support the central thesis that Vietnamese legal QA should
not use a single fixed retrieval strategy: dense retrieval suits direct
lookup, graph traversal suits amendment/replacement questions, hybrid
reasoning suits cases needing both, and clarification is required when the
target is underspecified. The key finding is that the second stage is not
merely a route corrector but a Chain-of-Thought step---the Two-stage Hybrid
surpasses even a perfect-label oracle. The clarification analysis adds that
surface-contextual ambiguity is handled reliably while semantic ambiguity
(missing-entity, multi-interpretation) remains open, motivating stronger
entity linking.

% ============================================================
\section{Limitations and Future Work}
\label{sec:limitations}
% ============================================================

Several limitations remain. (i)~Route labels are metadata-derived from
hop-count and cross-document structure rather than annotated by legal
experts. (ii)~Although paraphrasing mitigated template-based label leakage,
the benchmark---including the ambiguity and conversation stress tests---is
still synthetic and should be validated on a human-authored distribution.
(iii)~The strict split has no native \texttt{clarify} class, so
clarification depends entirely on Stage~2 and the ambiguity detector.
(iv)~Cross-document legal-effect relations (\textsc{AMENDS},
\textsc{REPLACES}) remain sparse, and the benchmark covers only
\emph{explicit} hybrid questions that name the documents to combine, not
\emph{implicit} ones. (v)~Stage~2 latency is $\approx\!119\times$ that of
Stage~1 routing and depends on local hardware; token usage is not logged.
(vi)~The hybrid source-ID mismatch means Hit@$k$ underrepresents true
hybrid quality, and human evaluation of legal correctness remains future
work.

Future directions include: (a)~a three-level cascade
XGBoost~$\rightarrow$~PhoBERT~$\rightarrow$~LLM to preserve routing
accuracy under the linguistic diversity of naturally phrased queries;
(b)~canonical \texttt{law\_id::article\_id} normalization with
cross-encoder re-ranking to address the retrieval format mismatch;
(c)~integrating manually reviewed \texttt{clarify} samples and implicit
hybrid questions into the training set; (d)~stronger legal relation
extraction for \textsc{AMENDS} and \textsc{REPLACES}; and (e)~human
evaluation of legal answer correctness and citation faithfulness.

% ============================================================
\section{Conclusion}
\label{sec:conclusion}
% ============================================================

We presented a Reasoning-aware Adaptive Router for Vietnamese Legal
Hybrid GraphRAG that combines a dense vector index of $68{,}835$ legal
text chunks, a Neo4j legal knowledge graph of $804{,}322$ nodes and
$1{,}968{,}621$ edges built from the Pháp~điển codification, an XGBoost
Stage~1 classifier, and a selective LLM Stage~2 verifier, dispatching each
query to dense retrieval, graph traversal, hybrid reasoning, or
clarification according to its reasoning demand and ambiguity level.

The central finding from the strict $600$-query benchmark is that
Stage~2 verification acts as a Chain-of-Thought mechanism. The
Two-stage Hybrid achieves Answer F1~$=0.661$, outperforming the
Single-stage Router (F1~$=0.530$) by $+0.131$ and, critically, the
Oracle Router with perfect routing labels (F1~$=0.606$) by $+0.055$.
This confirms that explicit Stage~2 reasoning over legal query intent
improves generation quality beyond what correct route selection alone
can provide. The Stage~1 XGBoost classifier separates the three routing
classes on the paraphrased test split with $0.995$ Macro-F1 at
millisecond latency, indicating that it learns robust structural features
rather than relying on surface templates.

On the ambiguity benchmarks, Stage~2 improves clarification F1 from
$0.000$ (Stage~1 alone) to $0.684$ (initial Stage~1+Stage~2) to
$0.840$ (Phase~3 full policy) with zero false positives throughout.
The conversation-aware \texttt{HistoryResolver} further improves the
separate $160$-query stress test to clarify F1~$=0.862$ and
history-resolution accuracy~$=0.850$.

Improving routing robustness through a three-level cascade,
normalizing canonical evidence IDs with cross-encoder re-ranking, and
validating benchmarks with human-reviewed legal queries remain the
primary directions for future investigation.

% ============================================================
\section*{Declarations}

\textbf{Funding.} This research received no specific grant from any
funding agency in the public, commercial, or not-for-profit sectors.

\textbf{Competing interests.} The authors declare that they have no
competing interests.

\textbf{Ethics approval.} Not applicable. This study does not involve
human participants or animals.

\textbf{Consent for publication.} Not applicable.

\textbf{Author contributions.} N.D.T.\ designed and implemented the
system, conducted all experiments, and drafted the manuscript.
P.V.H.\ and C.T.\ supervised the research, reviewed the experimental
design, and revised the manuscript. All authors read and approved
the final version.

\textbf{Data availability.} The Vietnamese legal QA routing dataset
and evaluation scripts are available from the corresponding author on
reasonable request.
% ============================================================

\bibliographystyle{spmpsci}
\bibliography{biblio}

\end{document}